\begin{derivation}[从\eqnrefs{eq:prescribed-current-linear-bath-hamiltonian-dp85}到\eqnrefs{eq:green-loss-cumulants-dp84}]
从正文相互作用绘景哈密顿量
\begin{equation*}
 H_I(t)
 =\sum_\lambda
 \left[f_\lambda(t)a_\lambda^\dagger
 +f_\lambda^*(t)a_\lambda\right]
\end{equation*}
出发。不同模式满足
\begin{equation*}
 [a_\lambda,a_{\lambda'}^\dagger]=\delta_{\lambda\lambda'},
 \qquad
 [a_\lambda,a_{\lambda'}]=0.
\end{equation*}
因此两个时刻的对易子可以逐项计算：
\begin{align*}
 [H_I(t),H_I(t')]
 &=\sum_{\lambda,\lambda'}
 \Bigl[
 f_\lambda(t)a_\lambda^\dagger+f_\lambda^*(t)a_\lambda,
 f_{\lambda'}(t')a_{\lambda'}^\dagger+f_{\lambda'}^*(t')a_{\lambda'}
 \Bigr]\\
 &=\sum_\lambda
 \left[
 f_\lambda^*(t)f_\lambda(t')
 -f_\lambda(t)f_\lambda^*(t')
 \right]\mathbb I.
\end{align*}
右端是 $c$ 数，因此
\begin{equation*}
 [H_I(t_1),[H_I(t_2),H_I(t_3)]]=0.
\end{equation*}
所以三阶及更高 Magnus 项逐项为零：
\begin{equation*}
 U_I=\exp(\Omega_1+\Omega_2),
\end{equation*}
其中
\begin{align*}
 \Omega_1
 &=-\frac{\ii}{\hbar}
 \int_{-\infty}^{\infty}\dd t\,H_I(t),\\
 \Omega_2
 &=-\frac1{2\hbar^2}
 \int_{-\infty}^{\infty}\dd t
 \int_{-\infty}^{t}\dd t'\,
 [H_I(t),H_I(t')].
\end{align*}
先算一阶项：
\begin{align*}
 \Omega_1
 &=-\frac{\ii}{\hbar}
 \sum_\lambda
 \left[
 a_\lambda^\dagger\int\dd t\,f_\lambda(t)
 +a_\lambda\int\dd t\,f_\lambda^*(t)
 \right]\\
 &=\sum_\lambda
 \left(\alpha_\lambda a_\lambda^\dagger
 -\alpha_\lambda^*a_\lambda\right),
\end{align*}
这里正是正文定义
\begin{equation*}
 \alpha_\lambda=-\frac{\ii}{\hbar}\int\dd t\,f_\lambda(t).
\end{equation*}
二阶项完全是纯虚 $c$ 数。把上面的对易子代入：
\begin{align*}
 \Omega_2
 &=-\frac1{2\hbar^2}
 \sum_\lambda
 \int\dd t\int_{-\infty}^{t}\dd t'\,
 \left[f_\lambda^*(t)f_\lambda(t')
 -f_\lambda(t)f_\lambda^*(t')\right]\\
 &=\ii\Phi,
\end{align*}
其中
\begin{equation*}
 \Phi
 =\frac1{\hbar^2}
 \sum_\lambda
 \Im\int\dd t\int_{-\infty}^{t}\dd t'\,
 f_\lambda(t)f_\lambda^*(t')
\end{equation*}
是实数。因此
\begin{equation*}
 U_I
 =\ee^{\ii\Phi}
 \exp\!\left[
 \sum_\lambda
 (\alpha_\lambda a_\lambda^\dagger
 -\alpha_\lambda^*a_\lambda)
 \right].
\end{equation*}
不同模式的位移生成元彼此对易：
\begin{equation*}
 [\alpha_\lambda a_\lambda^\dagger-\alpha_\lambda^*a_\lambda,
 \alpha_{\lambda'}a_{\lambda'}^\dagger-\alpha_{\lambda'}^*a_{\lambda'}]=0,
 \qquad \lambda\neq\lambda'.
\end{equation*}
所以普通指数可以精确分解：
\begin{equation*}
 U_I
 =\ee^{\ii\Phi}\prod_\lambda
 \exp(\alpha_\lambda a_\lambda^\dagger-\alpha_\lambda^*a_\lambda)
 =\ee^{\ii\Phi}\prod_\lambda D_\lambda(\alpha_\lambda).
\end{equation*}
作用在多模真空上，利用 $D(\alpha)|0\rangle=|\alpha\rangle$，得到正文\eqnrefs{eq:prescribed-current-multimode-coherent-state-dp85}。

接下来从相干态本身推导计数统计。Baker--Campbell--Hausdorff 恒等式给
\begin{equation*}
 D(\alpha)
 =\ee^{-|\alpha|^2/2}
 \ee^{\alpha a^\dagger}\ee^{-\alpha^*a}.
\end{equation*}
因为 $a|0\rangle=0$，
\begin{align*}
 |\alpha\rangle
 &=D(\alpha)|0\rangle\\
 &=\ee^{-|\alpha|^2/2}
 \sum_{n=0}^{\infty}\frac{\alpha^n}{n!}(a^\dagger)^n|0\rangle\\
 &=\ee^{-|\alpha|^2/2}
 \sum_{n=0}^{\infty}\frac{\alpha^n}{\sqrt{n!}}|n\rangle.
\end{align*}
因此单模量子数分布为
\begin{equation*}
 \Pr(N_\lambda=n)
 =|\langle n|\alpha_\lambda\rangle|^2
 =\ee^{-|\alpha_\lambda|^2}
 \frac{|\alpha_\lambda|^{2n}}{n!}.
\end{equation*}
其概率生成函数直接求和：
\begin{align*}
 G_\lambda(u)
 &\equiv\langle u^{N_\lambda}\rangle\\
 &=\ee^{-|\alpha_\lambda|^2}
 \sum_{n=0}^{\infty}
 \frac{(|\alpha_\lambda|^2u)^n}{n!}\\
 &=\exp\!\left[|\alpha_\lambda|^2(u-1)\right].
\end{align*}

环境总能量增量是
\begin{equation*}
 \Delta E_B=\sum_\lambda\hbar\omega_\lambda N_\lambda.
\end{equation*}
由于末态是不同模式相干态的直积，各 $N_\lambda$ 独立，所以总能量特征函数因子化：
\begin{align*}
 \chi_{\Delta E_B}(s)
 &=\left\langle
 \exp\!\left(\ii s\sum_\lambda\hbar\omega_\lambda N_\lambda\right)
 \right\rangle\\
 &=\prod_\lambda
 \left\langle \ee^{\ii s\hbar\omega_\lambda N_\lambda}\right\rangle\\
 &=\prod_\lambda
 \exp\!\left\{|\alpha_\lambda|^2
 [\ee^{\ii s\hbar\omega_\lambda}-1]\right\}.
\end{align*}
取对数后，乘积变成求和：
\begin{equation*}
 \ln\chi_{\Delta E_B}(s)
 =\sum_\lambda|\alpha_\lambda|^2
 [\ee^{\ii s\hbar\omega_\lambda}-1].
\end{equation*}
再使用正文定义的谱测度
\begin{equation*}
 \frac{\dd P_{\rm cl}}{\dd\omega}
 =\sum_\lambda|\alpha_\lambda|^2
 \delta(\omega-\omega_\lambda),
\end{equation*}
可逐项验证
\begin{align*}
 &\int_0^\infty\dd\omega\,
 \frac{\dd P_{\rm cl}}{\dd\omega}
 [\ee^{\ii s\hbar\omega}-1]\\
 &\qquad=\sum_\lambda|\alpha_\lambda|^2
 [\ee^{\ii s\hbar\omega_\lambda}-1].
\end{align*}
于是得到正文\eqnrefs{eq:green-loss-compound-poisson-cgf-dp84}。

最后，不直接引用“Poisson 的累积量都相同”，而从定义求导。第 $n$ 阶能量累积量是
\begin{equation*}
 \kappa_n
 =\left.\frac1{\ii^n}
 \frac{\dd^n}{\dd s^n}
 \ln\chi_{\Delta E_B}(s)\right|_{s=0}.
\end{equation*}
由于
\begin{equation*}
 \frac{\dd^n}{\dd s^n}
 \ee^{\ii s\hbar\omega}
 =(\ii\hbar\omega)^n \ee^{\ii s\hbar\omega},
\end{equation*}
在 $s=0$ 处得到
\begin{align*}
 \kappa_n(\Delta E_B)
 &=\int_0^\infty\dd\omega\,
 (\hbar\omega)^n
 \frac{\dd P_{\rm cl}}{\dd\omega},
\end{align*}
即正文\eqnrefs{eq:green-loss-cumulants-dp84}。$n=1,2$ 分别给均值和方差；这里的独立 Poisson 结构依赖“线性玻色环境 + 真空初态 + 给定经典轨迹”三个条件，其中任一条件失效都必须回到更一般的联合量子动力学。
\end{derivation}

